% =====================================================
% ROTEIRO PRÁTICO 04
% =====================================================

\section{Roteiro Prático 04 — Submissão de Artigo no SEPEI 2026}

% =====================================================
% OBJETIVOS
% =====================================================

\begin{boxObjetivos}
\textbf{Objetivos da Aula}

Ao final desta aula o estudante será capaz de:

\begin{itemize}
\item Compreender as regras e diretrizes de submissão do SEPEI 2026;
\item Adequar o trabalho desenvolvido ao template oficial do evento;
\item Realizar o cadastro na plataforma de submissão;
\item Efetuar a submissão formal do artigo científico.
\end{itemize}
\end{boxObjetivos}

% =====================================================
% INTRODUÇÃO
% =====================================================

\subsection{Introdução}

O Seminário de Ensino, Pesquisa, Extensão e Inovação do IFSC (SEPEI) é o principal evento científico da instituição. 

Nesta etapa do Projeto Integrador II, o artigo que vocês vêm desenvolvendo será submetido para avaliação e possível apresentação no SEPEI 2026.

As submissões estão abertas até \textbf{30 de abril de 2026} e podem ser para apresentações no formato presencial ou online.

% =====================================================
% PARTE 1
% =====================================================

\subsection{Parte 1 — Adequação ao Template Oficial}

\begin{boxAtividade}[Atividade 1 — Uso do Modelo do Evento]

O SEPEI 2026 exige rigor na formatação. Trabalhos fora do modelo serão \textbf{automaticamente desclassificados}.

Cada grupo deverá:

\begin{enumerate}
\item Acessar um dos modelos oficiais disponíveis:
    \begin{itemize}
        \item Modelo DOCX (Word): \url{https://ifsc.edu.br/documents/d/sepei/template-de-trabalho-docx-1}
        \item Modelo ODT (BrOffice): \url{https://ifsc.edu.br/documents/d/sepei/template-de-trabalho-odt-1}
        \item Modelo Google Docs: \url{https://docs.google.com/document/d/1tiXqEUAlOUpNYpoYjp1iFtrK-hS2mOKNLFNuarjzOqU/edit?usp=sharing}
    \end{itemize}
\item Transferir o conteúdo do artigo atual (desenvolvido nas aulas anteriores) para o modelo oficial do SEPEI.
\item Garantir que o trabalho tenha no \textbf{máximo 6 autores} (incluindo orientadores).
\end{enumerate}

\end{boxAtividade}

% =====================================================
% PARTE 2
% =====================================================

\subsection{Parte 2 — Estrutura e Seções do Artigo SEPEI}

\begin{boxAtividade}[Atividade 2 — Preenchimento das Seções do Template]

O documento oficial do SEPEI possui uma estrutura obrigatória que deve ser rigorosamente seguida (tamanho total de 10.000 a 12.000 caracteres com espaços). Abaixo estão as orientações para cada seção:

\begin{enumerate}
    \item \textbf{Cabeçalho (Título, Autores e Eixo):} Título em Arial 12, centralizado, negrito e caixa baixa. Máximo de 6 autores (com e-mail e vínculo em nota de rodapé). Escolha uma das 5 Divisões Temáticas (DT).\\
    \textit{Dica:} Preencha com atenção redobrada, pois os certificados serão emitidos com base nestes dados exatos.

    \item \textbf{Resumo e Palavras-chave:} Máximo de 250 palavras (sem parágrafo/citações), espaçamento simples, contendo introdução, objetivo, metodologia e resultados. 3 a 5 palavras-chave.\\
    \textit{Dica:} O resumo deve ser um texto em bloco único (1 parágrafo apenas), direto e objetivo.

    \item \textbf{Introdução:} Apresente a relevância do trabalho, contextualize o problema e apresente os objetivos, destacando a indissociabilidade entre ensino, pesquisa e extensão.\\
    \textit{Dica:} Espera-se de 2 a 3 parágrafos. Comece contextualizando o tema de forma ampla e finalize o último parágrafo declarando claramente o objetivo principal do projeto.

    \item \textbf{Fundamentação Teórica / Público Envolvido:} Descreva a base teórica (pesquisa/ensino) ou o público alvo (extensão).\\
    \textit{Dica:} Geralmente possui de 2 a 4 parágrafos. Utilize citações indiretas para embasar suas afirmações de forma sólida e evitar plágio.

    \item \textbf{Procedimentos Metodológicos:} Detalhe os materiais, métodos e etapas de desenvolvimento, alinhados com os objetivos.\\
    \textit{Dica:} Cerca de 1 a 3 parágrafos narrando o passo a passo da execução (o que foi feito, com quem, onde e como).

    \item \textbf{Resultados e Discussões:} Apresente os resultados obtidos e discuta-os. Relate os produtos ou registros gerados pelo projeto.\\
    \textit{Dica:} Costuma ser a seção mais longa (3 a 5 parágrafos). Use gráficos, tabelas ou figuras para ilustrar os resultados de forma clara, garantindo que todas as imagens sejam citadas no texto.

    \item \textbf{Considerações Finais:} Conclua se os objetivos foram alcançados. Destaque dificuldades, alcances e as contribuições para a formação.\\
    \textit{Dica:} De 1 a 2 parágrafos. Não cite novos autores aqui, foque em responder se o objetivo inicial foi cumprido e nos aprendizados adquiridos pelo grupo.

    \item \textbf{Referências e Fomento:} Siga o manual de normalização do IFSC. Se houver apoio financeiro, cite na seção correspondente.\\
    \textit{Dica:} Lembre-se de listar apenas as obras e os autores que foram efetivamente citados ao longo do texto.
\end{enumerate}

\end{boxAtividade}

% =====================================================
% PARTE 3
% =====================================================

\subsection{Parte 3 — Preparação para Submissão}

\begin{boxAtividade}[Atividade 3 — Tutorial de Submissão]

Antes de acessar o sistema para enviar o trabalho, é obrigatório assistir ao tutorial oficial do evento para evitar erros no preenchimento dos campos.

\begin{itemize}
\item Assista ao vídeo: \href{https://youtu.be/nArmca8ned4?si=ARA3-_HG5KX8TToX}{TUTORIAL - SUBMISSÕES DE TRABALHOS NO SEPEI}
\item Verifique se todos os metadados (resumo, título, autores) estão separados e fáceis de copiar.
\end{itemize}

\end{boxAtividade}

% =====================================================
% PARTE 4
% =====================================================

\subsection{Parte 4 — Cadastro e Submissão no Sistema}

\begin{boxAtividade}[Atividade 4 — Envio do Trabalho]

Com o documento pronto e revisado:

\begin{enumerate}
\item Acesse o sistema de submissões: \url{https://sepei.com.br/ojs/index.php/sepei/about/submissions}
\item Realize o \textbf{Cadastro} na plataforma (caso ainda não possua conta) ou faça o \textbf{Acesso}.
\item Inicie uma nova submissão.
\item Preencha todos os campos rigorosamente conforme as instruções do vídeo tutorial.
\item Marque as condições de submissão (ex: conformidade com o template e limite de autores).
\item Anexe o arquivo final do trabalho e conclua o envio.
\end{enumerate}

\end{boxAtividade}

% =====================================================
% CHECKLIST
% =====================================================

\subsection{Checklist Final da Aula}

\begin{boxAtividade}[Verificação Técnica]

Antes de finalizar:

\begin{itemize}
\item O trabalho foi transferido e formatado no template oficial do SEPEI?
\item O grupo conferiu o limite máximo de 6 autores?
\item Todos assistiram ao tutorial de submissão?
\item O cadastro na plataforma foi realizado?
\item A submissão foi concluída com sucesso e o comprovante salvo?
\end{itemize}

\end{boxAtividade}
